% !TEX program = xelatex
% !TEX options = -synctex=1 -interaction=nonstopmode -file-line-error
% example-neuroimaging.tex
% Demonstrates rossbeamerneuroimaging features: dark slides, \mnicoord,
% \colorbar, \brainslice, \contrastcard, atlas layout.
\documentclass[palette=brainlab]{rossbeamerneuroimaging}

\title{Neural Correlates of Face Processing}
\subtitle{An fMRI Study of the Fusiform Face Area}
\author{R.~Dunne, A.~Kanwisher}
\institute{Brain Imaging Centre, University of Manchester}
\date{June 2026}
\conference{Organisation for Human Brain Mapping, Seoul}

\begin{document}

% --- Title ---
\begin{frame}[plain]
  \titlepage
\end{frame}

% --- Dark content slide ---
\begin{frame}{Task Paradigm}
  Participants viewed alternating blocks of faces and houses in a
  block design (16\,s on, 16\,s off).

  \vskip8pt
  \begin{itemize}
    \item 24 healthy right-handed adults (12F, age $M = 27.4$)
    \item 3T Siemens Prisma, TR = 2000\,ms, 2\,mm isotropic
    \item \modalitybadge{fMRI} \quad \modalitybadge[rbBrainCool]{T1w}
  \end{itemize}

  \vskip10pt
  Peak activation for Faces~$>$~Houses at \mnicoord{42}{-52}{-18}
  in the right fusiform gyrus.
\end{frame}

% --- Contrast cards ---
\begin{frame}{Activation Results}
  \contrastcard{Faces $>$ Houses}{5.82}{186}{R Fusiform Gyrus (FFA)}
  \vskip6pt
  \contrastcard{Houses $>$ Faces}{4.31}{142}{L Parahippocampal Place Area}
  \vskip6pt
  \contrastcard{Faces $>$ Scrambled}{6.14}{224}{Bilateral STS / FFA}
  \vskip6pt
  \thresholdbadge{FWE $p < .05$, $k > 20$}
\end{frame}

% --- Colour bar example ---
\begin{frame}{Statistical Map}
  \begin{columns}[c]
    \begin{column}{0.75\textwidth}
      \centering
      \includegraphics[width=0.85\linewidth]{example-image}
    \end{column}
    \begin{column}{0.20\textwidth}
      \colorbar{$t$-value}{0}{8}
    \end{column}
  \end{columns}
  \vskip4pt
  \begin{center}
    \roiannotation{R FFA}\quad
    \roiannotation[rbBrainCool]{L PPA}\quad
    \roiannotation[rbBrainWarm]{STS}
  \end{center}
\end{frame}

% --- Brain slice ---
\begin{frame}{Axial Slices}
  \begin{columns}[c]
    \begin{column}{0.32\textwidth}
      \brainslice{example-image}{$z = -18$}
    \end{column}
    \begin{column}{0.32\textwidth}
      \brainslice{example-image}{$z = 0$}
    \end{column}
    \begin{column}{0.32\textwidth}
      \brainslice{example-image}{$z = 24$}
    \end{column}
  \end{columns}
\end{frame}

% --- Atlas layout ---
\begin{atlasframe}{Anatomical Reference}
  \atlasleft{example-image}{Sagittal view ($x = 42$)}
  \atlasright{%
    \atlaspanel{example-image}{Axial view ($z = -18$)}
    \atlaspanel{example-image}{Coronal view ($y = -52$)}
  }
\end{atlasframe}

% --- Coordinate table ---
\begin{frame}{Peak Coordinates}
  \centering
  \begin{coordtable}
    \coordrow{R Fusiform Gyrus}{42}{-52}{-18}{5.82}{186}
    \coordrow{L Parahippocampal}{-28}{-40}{-12}{4.31}{142}
    \coordrow{R Superior Temporal}{54}{-30}{8}{4.67}{98}
    \coordrow{L Inferior Frontal}{-44}{18}{24}{3.89}{64}
  \end{coordtable}
\end{frame}

% --- Pipeline ---
\begin{frame}{Processing Pipeline}
  \centering\vfill
  \pipelinestep{1}{Preprocessing}{Motion corr.,\\slice timing}
  \pipelinearrow
  \pipelinestep{2}{Normalisation}{MNI152,\\2\,mm}
  \pipelinearrow
  \pipelinestep{3}{Smoothing}{6\,mm FWHM\\Gaussian}
  \pipelinearrow
  \pipelinestep{4}{GLM}{SPM12,\\FWE corr.}
  \vfill
\end{frame}

\conclusionslide{Thank You}{r.dunne@manchester.ac.uk}

\end{document}
